\section{A uniform certificate for the actual finite residual support}
\label{sec:finite-range}

I now certify the support of the finite numerical oracle rather than
inferring that support from a mean, a variance or sampled extrema. The
certificate is conditional on the retained uniform arithmetic-error
certificate and the exact frozen model coefficients. Its finite quantile
support does not impose a bounded-support assumption on the underlying
continuous Gaussian financial model. The previously charged distributional
transfer error remains necessary when reporting a continuous-model price.
The convexity and moment inequalities used below are established results;
the project-specific construction is their uniform application to the
actual finite circuit and the resulting fair comparison between encodings.

\subsection{Exact latent support and the vertex maximum}

Let the emitted rounded marginal table, after removing its offset, consist
of integer values \(z_0,\ldots,z_{2^b-1}\), representing
\(2^{-q}z_j\). Its exact constructor is symmetric and increasing.
The independently checked maximum is therefore
\(v=2^{-q}z_{2^b-1}\), and every outer latent coordinate belongs to
\([-v,v]\). This is a statement about the actual integer table, whose
hash is retained. It is not a statement about the largest value of an
unbounded Gaussian random variable.

Write the frozen weights as \(w_i>0\), \(\sum_iw_i=1\), the outer
loadings as \(r_i\), and the log drifts as \(\mu_i\). Define
\(r_G=\sum_iw_i r_i\), \(\mu_G=\sum_iw_i\mu_i\),
\(s_i=r_i-r_G\), and \(d_i=\mu_i-\mu_G\). The ideal arithmetic to
geometric basket ratio at an exact rounded latent input \(y\) is
\[
 J(y)=\log\frac{A(y)}{G(y)}
     =\log\left(\sum_i w_i\exp(d_i+s_i\cdot y)\right).
                                                               \tag{R1}
\]
The shared inner factor cancels from this ratio. Since
\(\sum_iw_i d_i=0\) and \(\sum_iw_i s_i=0\), Jensen's inequality
for the exponential gives \(J(y)\geq0\).

For a direct convexity proof, put
\(\pi_i(y)=w_i\exp(d_i+s_i\cdot y)/\sum_jw_j\exp(d_j+s_j\cdot y)\).
Then
\[
 \nabla J=\sum_i\pi_i s_i,\qquad
 \nabla^2J=\sum_i\pi_i s_i s_i^{\mathsf T}
             -\left(\sum_i\pi_i s_i\right)
              \left(\sum_i\pi_i s_i\right)^{\mathsf T}.
\]
For every vector \(h\), the corresponding quadratic form is
\(\sum_i\pi_i(h\cdot s_i)^2-(\sum_i\pi_i h\cdot s_i)^2\),
a nonnegative variance. Thus \(J\) is convex. For \(v>0\), every
\(y\in[-v,v]^p\) is the convex combination of the vertices
\(v\sigma\), \(\sigma\in\{-1,1\}^p\), with coefficients
\[
 \lambda_\sigma(y)=\prod_{j=1}^p\frac{1+\sigma_j y_j/v}{2}.
\]
These coefficients are nonnegative, sum to one, and have weighted vertex
mean \(y\). Convexity therefore proves
\[
 0\leq J(y)\leq
 J_*:=\max_{\sigma\in\{-1,1\}^p}J(v\sigma).               \tag{R2}
\]
If \(v=0\), the box has one point and the same statement uses \(J(0)\).
The vertex search is exhaustive in the finite family \(p\leq7\);
it is not a numerical optimiser whose termination is assumed to find a
global maximum. The geometric outer log price has the independent bound
\[
 m_G(y)=\mu_G+r_G\cdot y
       \leq m_G^*:=\mu_G+v\sum_j\lvert r_{G,j}\rvert.        \tag{R3}
\]
The maxima in (R2) and (R3) need not occur at the same vertex. Multiplying
their separate nonnegative bounds below remains conservative.

\subsection{Method-specific ideal residual bounds}

Let \(K\) be the strike, \(H>0\) the payoff cap, \(h=H/K\), and
\(F_H(s)=\min((s-K)_+,H)\). The normalised capped payoff is one-Lipschitz
after scaling prices by \(H\). More precisely,
\[
 0\leq\frac{F_H(s)-F_H(t)}{H}\leq\frac{s-t}{H}
 \quad(s\geq t>0).
\]
As a function of log price, \(F_H(S_0e^m)/H\) has derivative
\(S_0e^m/H\) only while \(K<S_0e^m<K+H\), and derivative zero
outside that interval. Its global log-price Lipschitz constant is therefore
\(L_H=(K+H)/H=(1+h)/h\). This conclusion also follows by integrating
the bounded derivative across the two kinks.

For the direct method, the rounded inner table has exact support
\([-v_Z,v_Z]\), and the declared shared loading satisfies \(\lambda>0\).
The ideal rounded-grid normalised residual is bounded uniformly by
\[
 U_{\mathrm{direct}}=
 \min\left\{1,\ L_HJ_*,\
  \frac{S_0}{H}\exp(m_G^*+\lambda v_Z)(e^{J_*}-1)\right\}.   \tag{R4}
\]
Indeed, the arithmetic price equals the geometric price times \(e^{J(y)}\),
so its positive gap is at most the last expression. The other two terms
follow from the payoff cap and the log-price Lipschitz bound.

For a fixed strike sampled in the interval method, the residual is the
difference of two Gaussian distribution functions whose arguments differ
by \(J(y)/\lambda\). The Gaussian density is at most
\(\varphi(0)=(2\pi)^{-1/2}\). Hence
\[
 U_{\mathrm{interval}}=
 \min\{1,\ \varphi(0)J_*/\lambda\}.                       \tag{R5}
\]
This bound is uniform in the strike uniform, so the vertex search does not
need to enumerate that additional input.

For either fully conditioned method, averaging over the inner Gaussian
and over the capped payoff preserves the log-price slope bound \(L_H\)
and also gives the slope bound \(\varphi(0)/\lambda\).
The latter follows by writing the cap payoff as an integral of strike
indicators over \([K,K+H]\), dividing by \(H\), and differentiating
the bounded Gaussian distribution function. The positive price-gap bound
uses \(\mathbb E e^{\lambda Z}=e^{\lambda^2/2}\). Consequently,
\[
 U_{\mathrm{full}}=
 \min\left\{1,
   \min\{\varphi(0)/\lambda,L_H\}J_*,
   \frac{S_0}{H}\exp(m_G^*+\lambda^2/2)(e^{J_*}-1)\right\}.  \tag{R6}
\]
Equations (R4)--(R6) bound the ideal functions evaluated at the actual
rounded outer latent coordinates. They do not remove function-table
approximation, arithmetic truncation, or the continuous distributional
transfer allowance.

\subsection{Uniform body error, outward evaluation and lattice flooring}

Let \(a(y,\ldots)\) be the actual integer output and let the retained
arithmetic certificate prove a currency error no greater than \(\beta_A\).
Put \(\eta=\beta_A/H\). The same certificate remains conservative
when the exact rounded quantile values are taken as ideal body inputs:
their initial quantile error then becomes zero instead of the original
charged one-ULP bound. All retained domains still contain those inputs.
The ensuing error recurrences are monotone in the nonnegative preceding
error bounds, and their function domains shrink when those errors decrease.
Thus retaining the original \(\eta\) cannot undercharge the remaining
body arithmetic or interpolation error. In particular, this is a
pointwise body-error argument. An expectation-level distribution error is
not being converted to a pointwise or squared-error bound.

Use outward enclosures \(\overline J_*\), \(\overline L\) and
\(\overline U_{\mathrm{ideal}}\) in the appropriate preceding formula.
Then
\[
 0\leq X=2^{-q}a\leq
 U:=\min\{1,\overline U_{\mathrm{ideal}}+\eta+\epsilon_{\mathrm{num}}\},
 \qquad
 a\leq M:=\lfloor2^qU\rfloor.                            \tag{R7}
\]
The floor is justified by integrality: an integer no greater than a real
upper bound is no greater than its floor. It is not applied to an inward
rounding or to a sampled maximum. The conversion of the retained decimal
\(U\) to a rational number is exact. Therefore the final finite support
\(\{0,2^{-q},\ldots,M2^{-q}\}\) is exact conditional on the outward
range certificate, including the possible endpoint \(M=2^q\).

The generator evaluates at 220 decimal digits and adds a declared
\(10^{-100}\) outward allowance. Its finite numerical family has at
most eight assets, seven coordinates and 128 vertices, with exponential
arguments bounded in absolute value by 20. Correctly rounded Decimal
exponential and logarithm operations, rational-to-decimal conversions,
positive weighted sums and the bounded denominators in this family have
aggregate centre errors far smaller than that allowance. More explicitly,
the correctly rounded relative-error unit is at most
\(u=5\times10^{-220}\). Conversion of an exact exponent of magnitude
at most 20 contributes at most \(20u\) absolute error. Since
\(e^{20}<5\times10^8\), the mean-value bound for its exponential,
the eight weighted terms and their additions together give a conservative
absolute exponential-sum centre error below \(10^{-206}\).
The exact sum is at least one, so a perturbation of that size remains above
one half and the logarithm amplifies it by at most two. The subsequent
positive-gap candidate involves two exponentials with arguments of magnitude
at most 20 in the independently checked family and a coefficient
\(S_0/H\leq100\). Allowing all their conversions, products and
subtractions, the deliberately loose centre-error bound \(10^{-190}\)
still dominates those errors. The density slope divides only by
\(\lambda\geq0.05\), and its other denominator has \(h\geq0.01\),
so it does not invalidate that bound. Each added \(10^{-100}\)
allowance dominates this finite computation bound by ninety decimal orders.
Maximisation and minimisation select stored real enclosures and add no
new rounding operation. Weighted
arithmetic--geometric order keeps the exact exponential sum at least one,
so the logarithm does not amplify a small denominator. The independent
audit evaluates scaled log-sum-exp at 450 digits, uses a distinct positive
series for \(\pi\), and checks the retained upper enclosures against
all reference vertices with an explicit \(10^{-350}\) reference envelope.
That high-precision calculation supplements the numerical premise and
the uniform analytical proof; it is not a claim of a formally verified
implementation of Decimal itself.

To specify the density calculation's numerical premise in detail, the
generator evaluates Machin's identity
\(\pi=16\arctan(1/5)-4\arctan(1/239)\) by the alternating series
\(\arctan x=\sum_{n\geq0}(-1)^n x^{2n+1}/(2n+1)\).
Its loop includes the first power no greater than \(10^{-230}\).
For \(x=1/5\), it takes at most 166 updates; for \(x=1/239\),
at most 49. The margins at those powers are much larger than their
relative rounding errors, so these bounds also hold for the emitted
recurrences. For \(k\) successive rounded operations, the standard
product bound follows directly by expanding the product of their
\((1+\delta_j)\) factors:
\[
 \left\lvert\prod_{j=1}^k(1+\delta_j)-1\right\rvert
 \leq\frac{ku}{1-ku}=: \gamma_k\quad
 (\lvert\delta_j\rvert\leq u,\ ku<1).
\]
Each power update has at most two rounded multiplications. Since
\(\sum_{n\geq0}x^{2n+1}\leq x/(1-x^2)<0.21\) for both
arguments, the power, term-division and summation errors together are
bounded conservatively by
\(0.21(\gamma_{332}+\gamma_{167}+3u)<230u\).
The second argument's initial reciprocal error is at most \(u/239\);
the derivative of arctangent has magnitude at most one, so adding \(u\)
covers this perturbation. The alternating remainder is at most
\(10^{-230}\). Multiplying by the Machin coefficients and charging
their products and subtraction gives a pi error less than
\(20(231u+10^{-230})+100u<10^{-210}\).
The exact and computed pi therefore remain above three minus
\(10^{-210}\). The derivative of \((2\pi)^{-1/2}\) there is less
than 0.08 in magnitude. The correctly rounded reciprocal and square root
add less than \(10u\) absolute error. Thus the computed
\(\varphi(0)\) has absolute error below \(10^{-210}\).
Division by \(\lambda\geq0.05\) amplifies it by at most twenty,
which remains covered by the stated \(10^{-190}\) aggregate centre
bound before the \(10^{-100}\) outward pad.

For the other centre-error claims, exact rational exponent conversion has
absolute error at most \(20u\). The mean-value theorem and correct
rounding bound each exponential's error by
\(e^{20+20u}(20u)+u e^{20+20u}<6\times10^9u\).
Converting a weight in \([0,1]\), multiplying it by that exponential,
and summing at most eight positive terms contributes less than
\(10^{-206}\) in total. Dividing the logarithmic perturbation by a sum
at least one half and charging the rounded logarithm gives less than
\(3\times10^{-206}\). Every vertex obeys that bound, so taking the
maximum does too. The frozen weights have minimum at least 0.04. At least
one centred affine exponent is nonpositive because their weighted mean
is zero; consequently,
\(J_*\leq\log(0.96e^{20}+0.04)<19.96\).
Thus even its padded centre remains below 20. The independent audit also
checks the geometric-price exponential arguments in (R4) and (R6) against
the bound 20. Their conversions and products, two exponentials no greater
than \(5\times10^8\), and the coefficient at most 100 give an absolute
candidate centre error below \(10^{-190}\), including cancellation in
\(e^{\overline J_*}-1\): an absolute-error estimate is used, so a small
gap does not invalidate it. The slope product, final sum and bounded
normalisations satisfy the same loose bound. Every such candidate receives
an outward \(10^{-100}\) addition before the minimum is taken.
These estimates prove the numerical enclosure on the stated finite
parameter domain conditional on correct rounding of the declared Decimal
primitives; they do not extrapolate that envelope to arbitrary parameters.

\subsection{Exact transfer to mean and second-moment priors}

Let \(R=M2^{-q}\), and suppose previous certificates give
\(m\geq\mathbb E X\) and \(s\geq\mathbb E[X^2]\).
The support inequality \(X^2\leq RX\) gives the valid initial
tightenings
\[
 m_0=\min(m,R),\qquad s_0=\min(s,Rm_0,R^2).                \tag{R8}
\]
If \(M=0\), the finite output is identically zero and both bounds are
zero. Otherwise, write \(\delta=2^{-q}\). For every integer
\(a\in\{0,\ldots,M\}\) and every integer \(j\geq0\),
\((a-j)(a-j-1)\geq0\). For adjacent grid points
\(x=j\delta\), \(y=(j+1)\delta\), this proves the exact pointwise
linear moment certificate
\[
 X^2\geq(x+y)X-xy,
 \qquad \mathbb E X\leq\frac{s_0+xy}{x+y}.                \tag{R9}
\]
When \(s_0<R^2\), choose the unique \(j\in\{0,\ldots,M-1\}\)
with \(j^2\delta^2\leq s_0<(j+1)^2\delta^2\). The right side
of (R9) is the inverse of the lower convex hull of the finite squared-value
graph at \(s_0\). It is attained in the unrestricted lattice support
class by a mixture of \(x\) and \(y\) with second moment \(s_0\).
For \(j=0\), that mixture uses zero and \(\delta\), and (R9) reads
\(\mathbb E X\leq s_0/\delta\). If \(s_0\geq R^2\), the support
bound \(R\) is the corresponding hull bound. Denoting this inverse by
\(H_q(s_0,M)\), I retain
\[
 m_1=\min(m_0,H_q(s_0,M)),\qquad
 s_1=\min(s_0,Rm_1).                                      \tag{R10}
\]
Every step applies a pointwise inequality or a previously valid moment
bound, so it remains an upper bound for the actual finite financial law.
Sharpness is asserted only over all probability laws on the permitted
lattice, not over the narrower values attained by a particular financial
circuit. The circuit need not realise the extremising mixture. These
fair mean and second-moment priors are supplied to every eligible encoding,
including the single-flag range and clipping opponents.

The retained production-point checks evaluate the exact integer graph at
latent vertices, central and boundary addresses and reproducible random
inputs, then independently cross-check selected outputs using unbounded
Python integer replay. They can reveal a construction or coding failure,
but their maxima are not substituted for the uniform proof in (R7).
